\documentclass[11pt]{article}
\usepackage[margin=0.95in]{geometry}
\usepackage{amsmath,amssymb}
\usepackage{booktabs}
\usepackage{hyperref}
\newcommand{\code}[1]{\texttt{\small #1}}
\newcommand{\Exp}{\mathbb{E}}

\title{A semi-synthetic benchmark with a known exposure--response curve\\
\large Why the construction is what it is, and what it does and does not guarantee}
\author{}
\date{}

\begin{document}
\maketitle
\vspace{-2em}

\section{The problem}

We want to know whether the model recovers a causal exposure--response curve. On real data this
question has no answer: the true curve is exactly the unknown being estimated, so a recovered
curve can only be compared against another estimate, and any disagreement is uninterpretable.

The way out is to build data whose answer we chose. We keep the \emph{real} covariates --- real
PM$_{2.5}$, real gridMET, real census --- and synthesise only the \emph{outcome}, from a rate we
write down ourselves. The model trains on exactly the numbers that produced its target, and we can
ask whether the curve it recovers is the curve we injected.

This is worth doing only if two things hold. The benchmark must be \textbf{hard}: an unconfounded
exposure would let naive regression recover the truth, failing to separate a causal estimator from
a curve-fitter. And the ground truth must be \textbf{exact}: a Monte-Carlo ``truth'' carries its own
error bars, hiding small biases. Sections~\ref{sec:dgp} and~\ref{sec:truth} address these in turn.

\section{The generative model}\label{sec:dgp}

Let $t$ index days, $i$ index the $V$ ZCTAs. Write $X_{t,i}$ for the exposure (PM$_{2.5}$,
in $\mu$g/m$^3$) and $\tilde C^{(k)}_{t,i}$ for the $k$-th standardised confounder. The
per-capita Poisson rate is a sum of terms we specify:
\begin{equation}\label{eq:rate}
  \eta_{t,i} \;=\; \underbrace{r_0 + A_s \sin\!\Big(\tfrac{2\pi t}{365.25}\Big)
                  + \ell_{\mathrm{lat}}(i) + \ell_{\mathrm{lon}}(i)}_{\text{seasonal + geographic background}}
   \;+\; \beta\,X_{t,i} \;+\; \sum_k \gamma_k\, \tilde C^{(k)}_{t,i},
\end{equation}
\begin{equation}\label{eq:draw}
  \lambda_{t,i} \;=\; \max(\lambda_{\min},\,\eta_{t,i}), \qquad
  Y_{t,i} \;\sim\; \mathrm{Poisson}\big(\lambda_{t,i}\, o_i\big), \qquad
  o_i = \kappa \cdot \mathrm{pop}_i .
\end{equation}
Here $\beta$ \emph{is} the ground-truth dose--response and each $\gamma_k$ is a known confounder
effect. The floor $\lambda_{\min}$ keeps the rate positive; $o_i$ is the population offset, so
$\lambda$ is a rate per person and $\lambda o$ an expected count.

Three choices here are deliberate. The link is \textbf{identity, not log}, so $\Exp[Y]/o$ is
exactly linear in $X$: the ground truth is a closed form rather than an integral, and an ordinary
least-squares fit recovers $\beta$ and every $\gamma_k$ directly --- a check on the generator that
needs no model at all. The \textbf{exposure enters raw and confounders standardised}, so $\beta$ is
``per $\mu$g/m$^3$'' and comparable to the epidemiological literature, while each $\gamma_k$ is per
standard deviation and so comparable across covariates on very different scales. And the
dose--response is \textbf{same-day} --- $\lambda_{t,i}$ depends on $X_{t,i}$ and nothing else ---
the simplest defensible choice, with a sharp consequence for windowed forecasting
(Section~\ref{sec:window}).

\subsection*{Closing the confounding triangle}

Confounding needs two arms: the confounder must affect the outcome, \emph{and} it must affect the
exposure. Equation~\eqref{eq:rate} supplies the first through $\gamma_k$. The second is whatever
the real data happens to contain --- real income is correlated with real PM$_{2.5}$, but by an
amount we neither control nor know.

An optional generator supplies the missing arm on demand, decomposing each exposure into three
independent unit-variance parts,
\begin{equation}\label{eq:expo}
  X_e \;=\; m_e + s_e\Big(\sqrt{f^{\mathrm{dr}}}\, z\big(\textstyle\sum_k w_{ek}\tilde C^{(k)}\big)
        \;+\; \sqrt{f^{\mathrm{st}}}\, z\big(\mathcal{S}_{\rho,\phi}\varepsilon\big)
        \;+\; \sqrt{1 - f^{\mathrm{dr}} - f^{\mathrm{st}}}\, z(\varepsilon')\Big),
\end{equation}
where $z(\cdot)$ standardises and $\mathcal{S}_{\rho,\phi}$ imposes spatial and temporal
correlation. Because the parts are orthogonal by construction, $R^2$ of the exposure on the
confounders equals $f^{\mathrm{dr}}$ --- the confounding strength becomes a dial, not an accident.

The residual fraction $f^{\mathrm{no}} = 1 - f^{\mathrm{dr}} - f^{\mathrm{st}}$ is not slack: it is
the exposure variation remaining \emph{conditional on} the confounders. Driving it to zero destroys
positivity and makes the curve unidentifiable at any sample size, so the generator refuses it.

\begin{center}
\begin{tabular}{llc}
\toprule
Arm & Mechanism & Known? \\
\midrule
confounder $\to$ outcome  & $\gamma_k$ in \eqref{eq:rate}         & set in config \\
exposure $\to$ outcome    & $\beta$ in \eqref{eq:rate}            & set in config; the target \\
confounder $\to$ exposure & $f^{\mathrm{dr}}$ in \eqref{eq:expo}  & set in config \emph{when generated} \\
\bottomrule
\end{tabular}
\end{center}

Both worlds are supported. With \emph{real} exposures the third arm is uncontrolled but the data
are more realistic and the model needs no reconfiguration; with \emph{generated} exposures the
third arm is exact. The first is the default; one config switch selects the other.

\section{The ground truth}\label{sec:truth}

Two interventions are of interest. The \textbf{level} intervention $\mathrm{do}(X = a)$ sets the
exposure to a constant everywhere and traces the classical exposure--response curve. The
\textbf{shift} intervention $X \mapsto \delta X$ rescales the observed field, preserving its
spatial and temporal texture; $\delta = 1$ is the factual world. Under either, confounders are
left at their observed values --- that is precisely the g-computation operator.

The crucial point is that \eqref{eq:rate} is deterministic given the covariates. The Poisson draw
in \eqref{eq:draw} adds noise to the realised counts but not to their \emph{mean}, and the ground
truth is a statement about means. So we never simulate: we evaluate $\lambda$ under the
intervention directly, in closed form, including the floor. For the unfloored cells the slopes are
analytic,
\begin{equation}\label{eq:slopes}
  \frac{\partial}{\partial a}\,\Exp[Y \mid \mathrm{do}(X=a)] = \beta\, o,
  \qquad
  \frac{\partial}{\partial \delta}\,\Exp[Y \mid X \mapsto \delta X] = \beta\, o\, \Exp[X],
\end{equation}
so the two curves differ only by the exposure's own magnitude --- a cross-check we use to validate
them against each other. The floor is the only source of inexactness: it is a genuine
nonlinearity, so \eqref{eq:slopes} holds exactly only where $\eta > \lambda_{\min}$, which in
practice is about $99\%$ of cells.

\paragraph{Units are part of the estimand.} The data loader serves every covariate standardised.
Multiplying that channel by $\delta$ therefore implements $X \mapsto \mu + \delta(X - \mu)$, a
fan-out about the mean, not $X \mapsto \delta X$. The two are different interventions and can even
have opposite sign: if the year's mean exposure sits below the normalising mean $\mu$, ``increasing
by $10\%$'' in standardised units \emph{decreases} the expected outcome. Whichever convention the
evaluation uses, the ground truth must use the same one.

\section{The estimand is a sum over forecast windows}\label{sec:window}

The model is not a per-day regressor. From a start day $t_0$ it consumes $W$ daily input slices
($\tau = 0,\dots,W-1$, calendar day $t_0 + \tau$) and emits $H$ forecast leads off each
($k = 0,\dots,H-1$). All leads come from one forward pass --- a learned per-lead embedding on a
shared latent --- so the model is \emph{direct multi-horizon}, not autoregressive. Nothing is fed
back, so no error compounds and the closed form of Section~\ref{sec:truth} survives intact.

Prediction $(t_0,\tau,k)$ targets calendar day $s = t_0 + \tau + k$. Sweeping $t_0$ across the year
and summing gives the estimand
\begin{equation}\label{eq:estimand}
  G(\delta) \;=\; \frac{1}{V}\sum_{i}\sum_{t_0}\sum_{\tau=0}^{W-1}\sum_{k=0}^{H-1}
     f\big(\mathbf{X}^{(t_0)},\, A^{(t_0)}[\,m^\star \mapsto \delta A^{(t_0)}_{m^\star}]\big)_{\tau,k,i}\,.
\end{equation}
Two distinct multiplicities follow, and conflating them is the easiest way to misread a result.

First, a day in the interior of the year is summed $W \cdot H$ times, not $H$ times. A correction
of ``$\times H$ because the model predicts a window'' is short by a factor of $W$.

Second --- and this is the subtle one --- the intervention $A^{(t_0)}$ touches only that sample's
own $W$ input days, $t_0,\dots,t_0+W-1$. A prediction with $\tau + k \geq W$ therefore targets a
day whose exposure was left factual. Because the dose--response is same-day, those predictions
carry \emph{exactly zero} causal contrast, however large $\delta$ is. Writing $w_{\mathrm{all}}(s)$
and $w_{\mathrm{int}}(s)$ for the two counts,
\begin{equation}
  G(\delta) \;=\; \frac{1}{V}\sum_i o_i \sum_s
    \Big[ w_{\mathrm{int}}(s)\,\lambda_{s,i}(\delta) \;+\;
          \big(w_{\mathrm{all}}(s) - w_{\mathrm{int}}(s)\big)\,\lambda_{s,i}(1) \Big],
\end{equation}
with $w_{\mathrm{int}}/w_{\mathrm{all}} = \tfrac{1}{2}(W{+}1)/H$ in the interior (both ramp at the
year's edges). At $W = H = 30$ that is $465/900 \approx 52\%$; at $W = 3$, $H = 5$ it is
$6/15 = 40\%$. The remainder dilutes the curve without contributing signal.

The practical consequence: a perfectly calibrated model scored against a ground truth that
intervened on the \emph{whole} exposure field, rather than window-locally, appears to recover only
that fraction of the true effect. Both conventions are legitimate --- they are different estimands
--- but they must match on both sides, and $W$ and $H$ must be the values the model was trained
with.

\section{What this does and does not give you}

The construction guarantees that $\beta$ and every $\gamma_k$ are known exactly; that the ground
truth is computed from the same rate expression the generator draws from, so the two cannot drift;
and that a manifest of the resolved parameters is written beside every dataset, so the answer key
travels with the data rather than with a config file that may have moved on.

It does not make the benchmark realistic in every respect. The outcome is additive and same-day by
construction, so it rewards estimators matched to that structure. Poisson noise comes from a single
seeded stream, so residuals are correlated across diseases and years and should not be read as
signal. A confounder that is constant in the source data contributes nothing regardless of its
stated $\gamma$ --- harmless for identification, but not a variable the model must adjust for. And
the benchmark certifies that a correct adjustment recovers $\beta$; it says nothing about
robustness to confounders outside the specified set.

\end{document}
